\part{Phase 0: GPT-2 CPU-first correctness}

\chapter{GPT-2 architecture}
\label{ch:gpt2-architecture}
GPT-2 is the phase-0 architecture because it is simple enough to finish and complex enough to exercise the decoder runtime. \citep{radford2019gpt2}

\begin{table}[htbp]\centering\small
\begin{tabularx}{\textwidth}{L{0.22\textwidth}L{0.25\textwidth}Y}
\toprule
Module & Phase-0 form & Contract point \\
\midrule
Embeddings & token plus learned positions & No RoPE in GPT-2. \\
Attention & causal MHA & KV heads equal query heads. \\
MLP & GELU feed-forward & Biases present. \\
Cache & append-only & Full accumulated span visible. \\
Validation & PyTorch oracle & C++ and Rust compare within tolerance. \\
\bottomrule
\end{tabularx}
\caption{Phase-0 GPT-2 scope.}
\end{table}

\begin{intuitionbox}{Intuition}
The baseline proves the whole loop before newer variants arrive.
\end{intuitionbox}

\begin{definitionbox}{Mathematical or contract statement}
A GPT-2 block is pre-LayerNorm causal MHA plus GELU MLP with residual connections.
\end{definitionbox}

\begin{implementationbox}{Implementation interpretation}
Implement only the GPT-2 CPU float32 correctness path in phase 0.
\end{implementationbox}

\begin{verificationbox}{How to verify this works}
\begin{itemize}
\item Check shapes and dtypes at the boundary.
\item Compare deterministic outputs against PyTorch golden vectors.
\item Record tolerance, seed, model artifact, and backend.
\item Do not benchmark until validation passes.
\end{itemize}
\end{verificationbox}

\begin{warningbox}{Common failure modes}
\begin{itemize}
\item Letting framework defaults choose dtype or layout.
\item Testing only a batch size of one.
\item Depending on upstream checkpoint names in runtime code.
\item Treating benchmark output as validation.
\end{itemize}
\end{warningbox}

\begin{chapterexercises}
\item State the central invariant in one sentence.\\
\textit{Hint.} Use the conceptual guide vocabulary.
\item Construct a toy example with small shapes.\\
\textit{Hint.} Use $B=2$, $T=3$, $H=8$ where possible.
\item Name a validation test that would catch a plausible bug.\\
\textit{Hint.} Prefer generated golden vectors to handwritten long arrays.
\end{chapterexercises}

\chapter{Token and learned position embeddings}
\label{ch:gpt2-embeddings}
GPT-2 uses token embeddings and learned absolute position embeddings.



\begin{intuitionbox}{Intuition}
The first hidden tensor is where token identity meets position.
\end{intuitionbox}

\begin{definitionbox}{Mathematical or contract statement}
Let \(E^{tok}\in\R^{V\times H}\) be the token table and \(E^{pos}\in\R^{P\times H}\) be the learned position table. If \(c_b\) is the cache length before the current segment, then \(p_{b,t}=c_b+t\) and
\[
X_{b,t,:}=E^{tok}_{I_{b,t},:}+E^{pos}_{p_{b,t},:},
\qquad 0\leq p_{b,t}<P.
\]
\end{definitionbox}

\begin{implementationbox}{Implementation interpretation}
During decode, compute the position id from cache length; do not restart at zero.
\end{implementationbox}

\begin{verificationbox}{How to verify this works}
\begin{itemize}
\item Check shapes and dtypes at the boundary.
\item Compare deterministic outputs against PyTorch golden vectors.
\item Record tolerance, seed, model artifact, and backend.
\item Do not benchmark until validation passes.
\end{itemize}
\end{verificationbox}

\begin{warningbox}{Common failure modes}
\begin{itemize}
\item Letting framework defaults choose dtype or layout.
\item Testing only a batch size of one.
\item Depending on upstream checkpoint names in runtime code.
\item Treating benchmark output as validation.
\end{itemize}
\end{warningbox}

\begin{chapterexercises}
\item State the central invariant in one sentence.\\
\textit{Hint.} Use the conceptual guide vocabulary.
\item Construct a toy example with small shapes.\\
\textit{Hint.} Use $B=2$, $T=3$, $H=8$ where possible.
\item Name a validation test that would catch a plausible bug.\\
\textit{Hint.} Prefer generated golden vectors to handwritten long arrays.
\end{chapterexercises}

\chapter{LayerNorm}
\label{ch:layernorm}
LayerNorm normalises each token's hidden vector and applies learned scale and bias.

\[
\mu=\frac{1}{H}\sum_{j=1}^{H}x_j,\quad
\sigma^2=\frac{1}{H}\sum_{j=1}^{H}(x_j-\mu)^2,\quad
\operatorname{LN}(x)_j=\gamma_j\frac{x_j-\mu}{\sqrt{\sigma^2+\epsilon}}+\beta_j.
\]

\begin{intuitionbox}{Intuition}
LayerNorm is a per-token reduction over hidden channels.
\end{intuitionbox}

\begin{definitionbox}{Mathematical or contract statement}
For each token vector \(x\in\R^H\), the displayed formula is the phase-0 normalisation contract. The reduction axis is the hidden dimension only.
\end{definitionbox}

\begin{implementationbox}{Implementation interpretation}
Use config epsilon and float32 accumulation for the CPU reference.
\end{implementationbox}

\begin{verificationbox}{How to verify this works}
\begin{itemize}
\item Check shapes and dtypes at the boundary.
\item Compare deterministic outputs against PyTorch golden vectors.
\item Record tolerance, seed, model artifact, and backend.
\item Do not benchmark until validation passes.
\end{itemize}
\end{verificationbox}

\begin{warningbox}{Common failure modes}
\begin{itemize}
\item Letting framework defaults choose dtype or layout.
\item Testing only a batch size of one.
\item Depending on upstream checkpoint names in runtime code.
\item Treating benchmark output as validation.
\end{itemize}
\end{warningbox}

\begin{chapterexercises}
\item State the central invariant in one sentence.\\
\textit{Hint.} Use the conceptual guide vocabulary.
\item Construct a toy example with small shapes.\\
\textit{Hint.} Use $B=2$, $T=3$, $H=8$ where possible.
\item Name a validation test that would catch a plausible bug.\\
\textit{Hint.} Prefer generated golden vectors to handwritten long arrays.
\end{chapterexercises}

\chapter{Causal self-attention}
\label{ch:causal-self-attention}
Causal self-attention mixes visible previous tokens while hiding the future.

\[
\alpha_{b,h,t,s}
=
\operatorname{softmax}_{s}\left(
\frac{\langle Q_{b,h,t,:},K_{b,h,s,:}\rangle}{\sqrt{d_h}}
+M_{t,s}
\right),
\qquad
O_{b,h,t,:}=\sum_{s=0}^{S-1}\alpha_{b,h,t,s}V_{b,h,s,:}.
\]
For phase-0 GPT-2, \(H_{kv}=H_q\), so query head \(h\) reads key/value head \(h\). The causal mask \(M_{t,s}\) is \(0\) for visible key positions and \(-\infty\) for future positions.

\begin{intuitionbox}{Intuition}
Attention is content-addressed lookup under a no-future rule.
\end{intuitionbox}

\begin{definitionbox}{Mathematical or contract statement}
Scaled dot-product attention applies the causal mask before softmax. Masking after softmax is a different operation and generally changes probabilities.
\end{definitionbox}

\begin{implementationbox}{Implementation interpretation}
Split fused QKV projection into query, key, and value segments, then reshape by head.
\end{implementationbox}

\begin{verificationbox}{How to verify this works}
\begin{itemize}
\item Check shapes and dtypes at the boundary.
\item Compare deterministic outputs against PyTorch golden vectors.
\item Record tolerance, seed, model artifact, and backend.
\item Do not benchmark until validation passes.
\end{itemize}
\end{verificationbox}

\begin{warningbox}{Common failure modes}
\begin{itemize}
\item Letting framework defaults choose dtype or layout.
\item Testing only a batch size of one.
\item Depending on upstream checkpoint names in runtime code.
\item Treating benchmark output as validation.
\end{itemize}
\end{warningbox}

\begin{chapterexercises}
\item State the central invariant in one sentence.\\
\textit{Hint.} Use the conceptual guide vocabulary.
\item Construct a toy example with small shapes.\\
\textit{Hint.} Use $B=2$, $T=3$, $H=8$ where possible.
\item Name a validation test that would catch a plausible bug.\\
\textit{Hint.} Prefer generated golden vectors to handwritten long arrays.
\end{chapterexercises}

\chapter{GELU MLP}
\label{ch:gelu-mlp}
The GPT-2 MLP expands hidden states, applies GELU, then projects back.

\[
\operatorname{GELU}(x)\approx
\frac{x}{2}\left(1+\tanh\left(\sqrt{\frac{2}{\pi}}\left(x+0.044715x^3\right)\right)\right).
\]

\begin{intuitionbox}{Intuition}
The MLP gives each token nonlinear channel mixing.
\end{intuitionbox}

\begin{definitionbox}{Mathematical or contract statement}
For a token vector \(x\in\R^H\), \(W_1\in\R^{F\times H}\), \(b_1\in\R^F\), \(W_2\in\R^{H\times F}\), and \(b_2\in\R^H\),
\[
y=\operatorname{GELU}(xW_1^\top+b_1)W_2^\top+b_2\in\R^H.
\]
\end{definitionbox}

\begin{implementationbox}{Implementation interpretation}
Use the same GELU approximation as the PyTorch reference.
\end{implementationbox}

\begin{verificationbox}{How to verify this works}
\begin{itemize}
\item Check shapes and dtypes at the boundary.
\item Compare deterministic outputs against PyTorch golden vectors.
\item Record tolerance, seed, model artifact, and backend.
\item Do not benchmark until validation passes.
\end{itemize}
\end{verificationbox}

\begin{warningbox}{Common failure modes}
\begin{itemize}
\item Letting framework defaults choose dtype or layout.
\item Testing only a batch size of one.
\item Depending on upstream checkpoint names in runtime code.
\item Treating benchmark output as validation.
\end{itemize}
\end{warningbox}

\begin{chapterexercises}
\item State the central invariant in one sentence.\\
\textit{Hint.} Use the conceptual guide vocabulary.
\item Construct a toy example with small shapes.\\
\textit{Hint.} Use $B=2$, $T=3$, $H=8$ where possible.
\item Name a validation test that would catch a plausible bug.\\
\textit{Hint.} Prefer generated golden vectors to handwritten long arrays.
\end{chapterexercises}

\chapter{Transformer block composition}
\label{ch:block-composition}
A GPT-2 block composes LayerNorm, attention, residual addition, LayerNorm, MLP, and residual addition.



\begin{intuitionbox}{Intuition}
Residuals preserve an identity path through every block.
\end{intuitionbox}

\begin{definitionbox}{Mathematical or contract statement}
For block \(\ell\), pre-norm form is
\[
U^{(\ell)}=X^{(\ell)}+\operatorname{Attn}_{\ell}(\operatorname{LN}_{1,\ell}(X^{(\ell)})),
\qquad
X^{(\ell+1)}=U^{(\ell)}+\operatorname{MLP}_{\ell}(\operatorname{LN}_{2,\ell}(U^{(\ell)})).
\]
\end{definitionbox}

\begin{implementationbox}{Implementation interpretation}
Export block-boundary tensors so mismatches localise quickly.
\end{implementationbox}

\begin{verificationbox}{How to verify this works}
\begin{itemize}
\item Check shapes and dtypes at the boundary.
\item Compare deterministic outputs against PyTorch golden vectors.
\item Record tolerance, seed, model artifact, and backend.
\item Do not benchmark until validation passes.
\end{itemize}
\end{verificationbox}

\begin{warningbox}{Common failure modes}
\begin{itemize}
\item Letting framework defaults choose dtype or layout.
\item Testing only a batch size of one.
\item Depending on upstream checkpoint names in runtime code.
\item Treating benchmark output as validation.
\end{itemize}
\end{warningbox}

\begin{chapterexercises}
\item State the central invariant in one sentence.\\
\textit{Hint.} Use the conceptual guide vocabulary.
\item Construct a toy example with small shapes.\\
\textit{Hint.} Use $B=2$, $T=3$, $H=8$ where possible.
\item Name a validation test that would catch a plausible bug.\\
\textit{Hint.} Prefer generated golden vectors to handwritten long arrays.
\end{chapterexercises}

\chapter{Final normalisation and language-modelling head}
\label{ch:final-head}
The final LayerNorm and language-modelling head produce vocabulary logits.



\begin{intuitionbox}{Intuition}
The head turns hidden vectors into token scores.
\end{intuitionbox}

\begin{definitionbox}{Mathematical or contract statement}
If \(X\in\R^{B\times T\times H}\) and \(W_{lm}\in\R^{V\times H}\), logits satisfy \(Z_{b,t,:}=X_{b,t,:}W_{lm}^{\top}\in\R^V\). GPT-2 uses tied token embeddings in many checkpoints, but the internal artifact exposes the LM head explicitly.
\end{definitionbox}

\begin{implementationbox}{Implementation interpretation}
Load an explicit \tensorname{lm_head.weight} even when embeddings are tied.
\end{implementationbox}

\begin{verificationbox}{How to verify this works}
\begin{itemize}
\item Check shapes and dtypes at the boundary.
\item Compare deterministic outputs against PyTorch golden vectors.
\item Record tolerance, seed, model artifact, and backend.
\item Do not benchmark until validation passes.
\end{itemize}
\end{verificationbox}

\begin{warningbox}{Common failure modes}
\begin{itemize}
\item Letting framework defaults choose dtype or layout.
\item Testing only a batch size of one.
\item Depending on upstream checkpoint names in runtime code.
\item Treating benchmark output as validation.
\end{itemize}
\end{warningbox}

\begin{chapterexercises}
\item State the central invariant in one sentence.\\
\textit{Hint.} Use the conceptual guide vocabulary.
\item Construct a toy example with small shapes.\\
\textit{Hint.} Use $B=2$, $T=3$, $H=8$ where possible.
\item Name a validation test that would catch a plausible bug.\\
\textit{Hint.} Prefer generated golden vectors to handwritten long arrays.
\end{chapterexercises}

\chapter{PyTorch reference engine}
\label{ch:pytorch-reference-engine}
The PyTorch reference is the phase-0 oracle and vector exporter. \citep{pytorchDocs}

\begin{lstlisting}[language=Python,caption={Reference-forward sketch.}]
def forward(input_ids, cache=None):
    position_ids = make_position_ids(input_ids, cache)
    x = token_embedding(input_ids) + position_embedding(position_ids)
    for block in blocks:
        x = block(x, cache=cache)
    x = final_layer_norm(x)
    return x @ lm_head_weight.T
\end{lstlisting}

\begin{intuitionbox}{Intuition}
The reference should be boring, explicit, and inspectable.
\end{intuitionbox}

\begin{definitionbox}{Mathematical or contract statement}
PyTorch defines the expected \(f_\theta\) over converted artifacts.
\end{definitionbox}

\begin{implementationbox}{Implementation interpretation}
Construct the model from internal artifacts, not from an upstream pretrained class.
\end{implementationbox}

\begin{verificationbox}{How to verify this works}
\begin{itemize}
\item Check shapes and dtypes at the boundary.
\item Compare deterministic outputs against PyTorch golden vectors.
\item Record tolerance, seed, model artifact, and backend.
\item Do not benchmark until validation passes.
\end{itemize}
\end{verificationbox}

\begin{warningbox}{Common failure modes}
\begin{itemize}
\item Letting framework defaults choose dtype or layout.
\item Testing only a batch size of one.
\item Depending on upstream checkpoint names in runtime code.
\item Treating benchmark output as validation.
\end{itemize}
\end{warningbox}

\begin{chapterexercises}
\item State the central invariant in one sentence.\\
\textit{Hint.} Use the conceptual guide vocabulary.
\item Construct a toy example with small shapes.\\
\textit{Hint.} Use $B=2$, $T=3$, $H=8$ where possible.
\item Name a validation test that would catch a plausible bug.\\
\textit{Hint.} Prefer generated golden vectors to handwritten long arrays.
\end{chapterexercises}

\chapter{Exporting golden vectors}
\label{ch:golden-vectors}
Golden vectors are generated records used by all implementations.

\begin{lstlisting}[caption={Golden-vector sidecar sketch.}]
{
  "name": "smoke_hello",
  "model_family": "gpt2",
  "model_name": "gpt2",
  "input_ids": [15496, 995],
  "dtype": "float32",
  "tolerance": {"atol": 1e-4, "rtol": 1e-4}
}
\end{lstlisting}

\begin{intuitionbox}{Intuition}
A vector is a reproducible experiment, not a magic constant.
\end{intuitionbox}

\begin{definitionbox}{Mathematical or contract statement}
JSON carries metadata and NPZ carries larger arrays with a matching sidecar.
\end{definitionbox}

\begin{implementationbox}{Implementation interpretation}
Do not hard-code long vectors in LaTeX or source comments.
\end{implementationbox}

\begin{verificationbox}{How to verify this works}
\begin{itemize}
\item Check shapes and dtypes at the boundary.
\item Compare deterministic outputs against PyTorch golden vectors.
\item Record tolerance, seed, model artifact, and backend.
\item Do not benchmark until validation passes.
\end{itemize}
\end{verificationbox}

\begin{warningbox}{Common failure modes}
\begin{itemize}
\item Letting framework defaults choose dtype or layout.
\item Testing only a batch size of one.
\item Depending on upstream checkpoint names in runtime code.
\item Treating benchmark output as validation.
\end{itemize}
\end{warningbox}

\begin{chapterexercises}
\item State the central invariant in one sentence.\\
\textit{Hint.} Use the conceptual guide vocabulary.
\item Construct a toy example with small shapes.\\
\textit{Hint.} Use $B=2$, $T=3$, $H=8$ where possible.
\item Name a validation test that would catch a plausible bug.\\
\textit{Hint.} Prefer generated golden vectors to handwritten long arrays.
\end{chapterexercises}

\chapter{C++ correctness engine}
\label{ch:cpp-engine}
The C++ correctness engine is CPU float32 and consumes the same artifacts as PyTorch. \citep{cpp2020,cmakeDocs}

\begin{lstlisting}[language=C++,caption={C++ artifact view sketch.}]
struct TensorView {
    DType dtype;
    std::vector<std::size_t> shape;
    std::span<const std::byte> bytes;
};
\end{lstlisting}

\begin{intuitionbox}{Intuition}
C++ earns trust by explicit memory ownership and checked views.
\end{intuitionbox}

\begin{definitionbox}{Mathematical or contract statement}
The C++ engine implements config, loader, tensor views, CPU ops, modules, cache, validation, and CLI.
\end{definitionbox}

\begin{implementationbox}{Implementation interpretation}
Add optimisation only after scalar correctness passes.
\end{implementationbox}

\begin{verificationbox}{How to verify this works}
\begin{itemize}
\item Check shapes and dtypes at the boundary.
\item Compare deterministic outputs against PyTorch golden vectors.
\item Record tolerance, seed, model artifact, and backend.
\item Do not benchmark until validation passes.
\end{itemize}
\end{verificationbox}

\begin{warningbox}{Common failure modes}
\begin{itemize}
\item Letting framework defaults choose dtype or layout.
\item Testing only a batch size of one.
\item Depending on upstream checkpoint names in runtime code.
\item Treating benchmark output as validation.
\end{itemize}
\end{warningbox}

\begin{chapterexercises}
\item State the central invariant in one sentence.\\
\textit{Hint.} Use the conceptual guide vocabulary.
\item Construct a toy example with small shapes.\\
\textit{Hint.} Use $B=2$, $T=3$, $H=8$ where possible.
\item Name a validation test that would catch a plausible bug.\\
\textit{Hint.} Prefer generated golden vectors to handwritten long arrays.
\end{chapterexercises}

\chapter{Rust correctness engine}
\label{ch:rust-engine}
The Rust correctness engine mirrors the C++ semantics while using ownership and typed errors. \citep{rustBook,rust2024edition}

\begin{lstlisting}[caption={Rust loader trait sketch.}]
pub trait WeightLoader {
    fn tensor<'a>(&'a self, name: &str) -> Result<Tensor<'a>, Error>;
}
\end{lstlisting}

\begin{intuitionbox}{Intuition}
Rust helps prevent stale or aliased cache and tensor views.
\end{intuitionbox}

\begin{definitionbox}{Mathematical or contract statement}
Rust consumes the same artifacts and compares against the same vectors.
\end{definitionbox}

\begin{implementationbox}{Implementation interpretation}
Use \code{Result} for malformed artifacts and validation failures.
\end{implementationbox}

\begin{verificationbox}{How to verify this works}
\begin{itemize}
\item Check shapes and dtypes at the boundary.
\item Compare deterministic outputs against PyTorch golden vectors.
\item Record tolerance, seed, model artifact, and backend.
\item Do not benchmark until validation passes.
\end{itemize}
\end{verificationbox}

\begin{warningbox}{Common failure modes}
\begin{itemize}
\item Letting framework defaults choose dtype or layout.
\item Testing only a batch size of one.
\item Depending on upstream checkpoint names in runtime code.
\item Treating benchmark output as validation.
\end{itemize}
\end{warningbox}

\begin{chapterexercises}
\item State the central invariant in one sentence.\\
\textit{Hint.} Use the conceptual guide vocabulary.
\item Construct a toy example with small shapes.\\
\textit{Hint.} Use $B=2$, $T=3$, $H=8$ where possible.
\item Name a validation test that would catch a plausible bug.\\
\textit{Hint.} Prefer generated golden vectors to handwritten long arrays.
\end{chapterexercises}

\chapter{Cross-language validation}
\label{ch:cross-language}
Cross-language validation compares PyTorch, C++, and Rust over one artifact directory.



\begin{intuitionbox}{Intuition}
Three runtimes are valuable only if disagreement is visible.
\end{intuitionbox}

\begin{definitionbox}{Mathematical or contract statement}
Shapes, dtype, hidden states, logits, greedy ids, generation, and cache equivalence are required checks.
\end{definitionbox}

\begin{implementationbox}{Implementation interpretation}
Build one comparison policy and use it for every implementation.
\end{implementationbox}

\begin{verificationbox}{How to verify this works}
\begin{itemize}
\item Check shapes and dtypes at the boundary.
\item Compare deterministic outputs against PyTorch golden vectors.
\item Record tolerance, seed, model artifact, and backend.
\item Do not benchmark until validation passes.
\end{itemize}
\end{verificationbox}

\begin{warningbox}{Common failure modes}
\begin{itemize}
\item Letting framework defaults choose dtype or layout.
\item Testing only a batch size of one.
\item Depending on upstream checkpoint names in runtime code.
\item Treating benchmark output as validation.
\end{itemize}
\end{warningbox}

\begin{chapterexercises}
\item State the central invariant in one sentence.\\
\textit{Hint.} Use the conceptual guide vocabulary.
\item Construct a toy example with small shapes.\\
\textit{Hint.} Use $B=2$, $T=3$, $H=8$ where possible.
\item Name a validation test that would catch a plausible bug.\\
\textit{Hint.} Prefer generated golden vectors to handwritten long arrays.
\end{chapterexercises}

\chapter{Cache-equivalence testing}
\label{ch:cache-equivalence}
Cache equivalence is mandatory in phase 0.



\begin{intuitionbox}{Intuition}
The cache is correct only if removing it does not change the answer.
\end{intuitionbox}

\begin{definitionbox}{Mathematical or contract statement}
For a prompt \(u_{1:m}\) and next token \(v\), the final-position logits from full recomputation on \((u_{1:m},v)\) must match the logits from prefill on \(u_{1:m}\) followed by cached decode on \(v\), within tolerance. The comparison also checks that every layer's cache length increases from \(m\) to \(m+1\).
\end{definitionbox}

\begin{implementationbox}{Implementation interpretation}
Test final-position logits, cache lengths, and per-layer append behaviour.
\end{implementationbox}

\begin{verificationbox}{How to verify this works}
\begin{itemize}
\item Check shapes and dtypes at the boundary.
\item Compare deterministic outputs against PyTorch golden vectors.
\item Record tolerance, seed, model artifact, and backend.
\item Do not benchmark until validation passes.
\end{itemize}
\end{verificationbox}

\begin{warningbox}{Common failure modes}
\begin{itemize}
\item Letting framework defaults choose dtype or layout.
\item Testing only a batch size of one.
\item Depending on upstream checkpoint names in runtime code.
\item Treating benchmark output as validation.
\end{itemize}
\end{warningbox}

\begin{chapterexercises}
\item State the central invariant in one sentence.\\
\textit{Hint.} Use the conceptual guide vocabulary.
\item Construct a toy example with small shapes.\\
\textit{Hint.} Use $B=2$, $T=3$, $H=8$ where possible.
\item Name a validation test that would catch a plausible bug.\\
\textit{Hint.} Prefer generated golden vectors to handwritten long arrays.
\end{chapterexercises}

\chapter{Baseline benchmarking}
\label{ch:baseline-benchmarking}
Benchmarking comes after correctness and measures prefill latency, decode latency, end-to-end generation, and throughput.



\begin{intuitionbox}{Intuition}
A benchmark is a measurement of a known-correct path.
\end{intuitionbox}

\begin{definitionbox}{Mathematical or contract statement}
Useful metrics include prompt length, generated tokens, batch size, latency, tokens per second, and cache bytes.
\end{definitionbox}

\begin{implementationbox}{Implementation interpretation}
Store benchmark suites separately from golden vectors.
\end{implementationbox}

\begin{verificationbox}{How to verify this works}
\begin{itemize}
\item Check shapes and dtypes at the boundary.
\item Compare deterministic outputs against PyTorch golden vectors.
\item Record tolerance, seed, model artifact, and backend.
\item Do not benchmark until validation passes.
\end{itemize}
\end{verificationbox}

\begin{warningbox}{Common failure modes}
\begin{itemize}
\item Letting framework defaults choose dtype or layout.
\item Testing only a batch size of one.
\item Depending on upstream checkpoint names in runtime code.
\item Treating benchmark output as validation.
\end{itemize}
\end{warningbox}

\begin{chapterexercises}
\item State the central invariant in one sentence.\\
\textit{Hint.} Use the conceptual guide vocabulary.
\item Construct a toy example with small shapes.\\
\textit{Hint.} Use $B=2$, $T=3$, $H=8$ where possible.
\item Name a validation test that would catch a plausible bug.\\
\textit{Hint.} Prefer generated golden vectors to handwritten long arrays.
\end{chapterexercises}
